\documentclass[../../../../main.tex]{subfiles}

\begin{document}

\section*{第1問}

\begin{proof}[解]
$A(1,0)$ とする。$\overline{AP}$ の中点 $M_1$ とする。
$\angle M_1OA=\theta$ とすると，$\overline{OA}$ から $\overline{OQ}$ へ向かう角は $4\theta$ であるから，$Q(\cos4\theta,\sin4\theta)$ となる。 $\cdots$①
ここで右下の三角形 $\triangle M_1OA$ に注目して，
\[
\sin\theta=\sqrt{1-r^2}, \quad \cos\theta=r \quad \cdots \text{②}
\]
だから，
\[
\cos4\theta=2\cos^22\theta-1=2(2\cos^2\theta-1)^2-1=2(2r^2-1)^2-1 \quad (\because②) = 8r^4-8r^2+1
\]
\[
\sin4\theta=2\sin2\theta\cos2\theta=4\sin\theta\cos\theta\cdot(2\cos^2\theta-1)=4r\sqrt{1-r^2}(2r^2-1)
\]
となるので，①から
\[
Q\left(8r^4-8r^2+1,\ 4r(2r^2-1)\sqrt{1-r^2}\right)
\]
\end{proof}

\newpage

\section*{第2問}

\begin{proof}[解]
$A=\{20,21,22,23,24\}$，$B=\{25,26,27,28,29\}$，$C=\{80,81,82,83,84\}$，
$D=\{85,86,87,88,89\}$ とする。題意の試行では，$A,B$ のうちから1つ，$C,D$ のうちから1つ数をとり出す。その組は以下で全てである。

\begin{center}
\begin{tabular}{c|c|c}
 & $A$ & $B$ \\
\hline
$C$ & ア & イ \\
\hline
$D$ & ウ & エ
\end{tabular}
\end{center}

このうち，アでは必ず $S\ge S'$，エでは必ず $S<S'$ となる。 $\cdots$①

そこで，以下イ，ウについて考える。

\textbf{$1^\circ$ イの時}

$B,C$ の数は，四捨五入すると各々30,80になるから，$S'=2400$ である。一方，$S$ は以下のようになる。（$\bigcirc$ は $S<S'$ を，$\times$ は $S\ge S'$ を表す）

\begin{center}
\begin{tabular}{c|c|c|c|c|c}
$C\backslash B$ & 25 & 26 & 27 & 28 & 29 \\
\hline
80 & $\bigcirc$ & $\bigcirc$ & $\bigcirc$ & $\bigcirc$ & $\bigcirc$ \\
\hline
81 & $\bigcirc$ & $\bigcirc$ & $\bigcirc$ & $\bigcirc$ & $\bigcirc$ \\
\hline
82 & $\bigcirc$ & $\bigcirc$ & $\bigcirc$ & $\bigcirc$ & $\bigcirc$ \\
\hline
83 & $\bigcirc$ & $\bigcirc$ & $\bigcirc$ & $\bigcirc$ & $\times$ \\
\hline
84 & $\bigcirc$ & $\bigcirc$ & $\bigcirc$ & $\bigcirc$ & $\times$
\end{tabular}
\end{center}

\textbf{$2^\circ$ ウの時}

$A,D$ の数は，四捨五入すると各々20,90になるから，$S'=1800$ である。$S$ の表は以下

\begin{center}
\begin{tabular}{c|c|c|c|c|c}
$D\backslash A$ & 20 & 21 & 22 & 23 & 24 \\
\hline
85 & $\bigcirc$ & $\bigcirc$ & $\times$ & $\times$ & $\times$ \\
\hline
86 & $\bigcirc$ & $\times$ & $\times$ & $\times$ & $\times$ \\
\hline
87 & $\bigcirc$ & $\times$ & $\times$ & $\times$ & $\times$ \\
\hline
88 & $\bigcirc$ & $\times$ & $\times$ & $\times$ & $\times$ \\
\hline
89 & $\bigcirc$ & $\times$ & $\times$ & $\times$ & $\times$
\end{tabular}
\end{center}

①，$1^\circ$，$2^\circ$ から，全てのえらび方 $10^2=100$ 通りのうち，$S<S'$ となるのは，
\[
25+23+6=54 \text{（通り）}
\]
（① ，$1^\circ$ ，$2^\circ$ の順に対応）

だから，もとめるカクリツは $\dfrac{54}{100}=\dfrac{27}{50}$
\end{proof}

\newpage

\section*{第3問}

\begin{proof}[解]
$t>0$ に対し，$x=t$ とすると，$P(t,t^2)$ である。$P$ における接線 $\ell$ は $(x^2)'=2x$ より，
\[
\ell: y=2tx-t^2 \quad \left(\text{方向ベクトル}\ \vec\ell=\begin{pmatrix}1\\2t\end{pmatrix}\right)
\]
である。この時，直線 $AP$ の方程式は，
\[
AP: y=\frac{t^2}{t+\frac13}\left(x+\frac13\right) \quad \left(\text{方向ベクトル}\ \vec{m}=\begin{pmatrix}t+\frac13\\t^2\end{pmatrix}\right)
\]
となる。$y$軸と平行なベクトル $\vec n=\begin{pmatrix}0\\1\end{pmatrix}$ をおき，$\vec\ell$ と $\vec m$，$\vec\ell$ と $\vec n$ のなす角を各々 $\theta_1,\theta_2$ とすると，（$0\le\theta_1,\theta_2\le\pi$）$t>0$ から $\theta_1,\theta_2\ne\pi/2$ であるから，題意から，
\[
\tan\theta_2=\pm\tan\theta_1 \quad \cdots \text{①}
\]
となる。
\[
\tan\theta_1=\frac{|t^2-2t(t+\frac13)|}{\vec\ell\cdot\vec m}=\frac{|t^2+\frac23t|}{t+\frac13+2t^3}=\frac{t^2+\frac23t}{2t^3+t+\frac13} \quad (\because t>0)
\]
\[
\tan\theta_2=\frac{|1|}{\vec\ell\cdot\vec n}=\frac{1}{2t}
\]
だから，$\tan\theta_1,\tan\theta_2>0$ であることより，①で複号正を採用して代入して
\[
\frac{1}{2t}=\frac{t^2+\frac23t}{2t^3+t+\frac13}
\]
\[
\therefore\ 2t^3+t+\frac13=2t^3+\frac43t^2 \quad\therefore\ t=1\ (\because t>0)
\]
となり，$P(1,1)$ である。

\medskip
\noindent\textbf{[解2]}

$\ell,AP$ の傾きを各々 $\ell=\tan\alpha,\ m=\tan\beta$ とおく。（$-\pi/2<\alpha,\beta<\pi/2$）
\[
\ell=2t,\quad m=\frac{t^2}{t+\frac13}
\]
である。右図のように $\theta_1,\theta_2$ をとると，どのような場合にも，
\[
\begin{cases}
\theta_1=\pi-(\alpha-\beta) \\
\theta_2=\dfrac{\pi}{2}-\alpha
\end{cases}
\quad \cdots \star
\]
となる。題意から
\[
\tan\theta_2=\pm\tan\theta_1
\]
であり，$\star$から，
\[
\frac{1}{\tan\alpha}=\pm\tan(\alpha-\beta)=\pm\frac{2t-\frac{t^2}{t+1/3}}{1+2t\cdot\frac{t^2}{t+1/3}}=\frac{t^2+\frac23t}{2t^3+t+\frac13}
\]
\[
\therefore\ \frac{1}{2t}=\frac{t^2+\frac23t}{2t^3+t+\frac13}
\]
（以下同様）
\end{proof}

\newpage

\section*{第4問}

\begin{proof}[解]
$m,n\in\mathbb{Z}_{>0}$，$m<n$，$0<x<1$ $\cdots$①である。$f(t)=\left(1+\dfrac{x}{t^2}\right)^t$ とする。$f(t)>0$ のゆえ，自然対数をとってから微分して，
\[
\frac{f'(t)}{f(t)}=\log\left(1+\frac{x}{t^2}\right)+\frac{t}{1+x/t^2}(-2)\frac{1}{t^3} = \log\left(1+\frac{x}{t^2}\right)-\frac{2}{x+t^2} \quad \cdots \text{②}
\]
この右辺を $g(t)$ とする。
\[
g'(t)=\frac{1}{1+x/t^2}(-2)\frac{1}{t^3}+2\left(\frac{1}{t^2+x}\right)^2\cdot2t = -\frac2t\cdot\frac{1}{x+t^2}+4t\left(\frac{1}{x+t^2}\right)^2 = \frac{2(t^2-x)}{t(x+t^2)^2}
\]
だから，①より，$1\le t$ の時，$g'(t)\ge0$ で，$g(t)$ は単調増加となる。これと
\[
g(t)\xrightarrow{t\to+\infty}0
\]
から，$1\le t$ の時，$g(t)<0$ となる。したがって，②から，$1\le t$ の時，$f'(t)<0$ つまり $f(t)$ は単調減少である。 $\cdots$③

①③から，$f(n)<f(m)$ つまり，
\[
\left(1+\frac{x}{m^2}\right)^m>\left(1+\frac{x}{n^2}\right)^n
\]
である。
\end{proof}

\newpage

\section*{第5問}

\begin{proof}[解]
$f(x)$ は任意の $x$ に対し，$f(x)=f(x+2\pi)$ をみたす。 $\cdots$①

\begin{enumerate}
  \item[(1)] $s=t-x$ とすると，
  \[
  \int_0^{2\pi}f(t-x)\sin t\,dt=\int_{-x}^{2\pi-x}f(s)\sin(x+s)ds
  \]
  \[
  =\int_{-x}^{0}f(s)\sin(x+s)ds+\int_0^{2\pi-x}f(s)\sin(x+s)ds
  \]
  \[
  =\int_{2\pi-x}^{2\pi}f(p-2\pi)\sin(x+p-2\pi)dp+\int_0^{2\pi-x}f(s)\sin(x+s)ds \quad (\text{第1項で}\ p=s+2\pi\text{とした})
  \]
  \[
  =\int_{2\pi-x}^{2\pi}f(p)\sin(x+p)dp+\int_0^{2\pi-x}f(s)\sin(x+s)ds \quad (\because①，及び\sin(p+2\pi)=\sin p)
  \]
  \[
  =\int_0^{2\pi}f(s)\sin(x+s)ds
  \]

  \item[(2)] (1)から，
  \[
  \int_0^{2\pi}f(t)\sin(x+t)dt=cf(x) \quad \cdots \text{②}
  \]
  が $x$ についての恒等式になれば良い。以下 $S=\sin x,\ P=\cos x$ とする。②の左辺を変形する。
  \[
  (\text{②の左辺})=\int_0^{2\pi}\{f(t)\cos t\cdot S+f(t)\sin t\cdot P\}dt = S\int_0^{2\pi}f(t)\cos t\,dt+P\int_0^{2\pi}f(t)\sin t\,dt \quad \cdots \text{③}
  \]
  $A=\displaystyle\int_0^{2\pi}f(t)\cos t\,dt,\ B=\int_0^{2\pi}f(t)\sin t\,dt \quad \cdots \text{④}$ とおいて，②③より，
  \[
  f(x)=\frac1c\{AS+BP\}=\frac{A}{c}S+\frac{B}{c}P \quad \cdots \text{⑤}
  \]
  ⑤に代入して，
  \[
  A=\int_0^{2\pi}\left\{\frac{A}{c}SP+\frac{B}{c}P^2\right\}dx=\left[-\frac{A}{4c}\cos2x+\frac{B}{2c}\left(x+\frac12\sin2x\right)\right]_0^{2\pi}=\frac{B}{c}\pi \quad \cdots \text{⑥}
  \]
  \[
  B=\int_0^{2\pi}\left\{\frac{A}{c}S^2+\frac{B}{c}SP\right\}dx=\left[\frac{A}{2c}\left(x-\frac12\sin2x\right)-\frac{B}{4c}\cos2x\right]_0^{2\pi}=\frac{A}{c}\pi \quad \cdots \text{⑦}
  \]
  ⑤⑥及び $f(0)=1$ から $(A,B)=(0,0)$ ではなく，$c=B$ である（$\because$④）ことから，
  \[
  (A,B,c)=(\pi,\pi,\pi)
  \]
  となり，④から
  \[
  f(x)=\sin x+\cos x, \quad c=\pi
  \]
  を得る。
\end{enumerate}
\end{proof}

\newpage

\section*{第6問}

\begin{flushright}
\footnotesize\textit{（(a)(b)いずれか一方を選択する問題。(a)のみ解答あり。(b)は原稿は空欄）}
\end{flushright}

\begin{proof}[解]
任意の $n\in\mathbb{N}$ に対し，$|a_n|\le1$ となることを示す。まず，$a_N=1$ 及び $a_{n+1}=2a_n^2-1$ から，$n\ge N$ の時，$a_n=1$ で成立する。そこで以下 $n\le N$ の時を考える。$n=k$（$2\le k\le N$）で成立を仮定すると，
\[
a_{k-1}^2=\frac12(a_k+1)
\]
\[
\therefore\ 0\le a_{k-1}^2\le1 \quad (\because|a_k|\le1)
\]
\[
\therefore\ |a_{k-1}|\le1
\]
となり $n=k-1$ でも成立。したがって，帰納法により，任意の $n\in\mathbb{N}$ に対して $|a_n|\le1$ となる。したがって，(1)の主張 $|a_1|\le1$ が従う。

この時，$a_n=\cos\theta_n$ とおける。漸化式から
\[
\cos\theta_{n+1}=\cos2\theta_n \quad\therefore\ \theta_{n+1}=2\theta_n \quad\therefore\ \theta_n=2^{n-1}\theta_1 \quad \cdots \text{①}
\]
これと $\theta_N=2k\pi$（$k\in\mathbb{Z}$）から，
\[
\theta_1=\frac{\theta_N}{2^{N-1}}=\frac{k\pi}{2^{N-2}}
\]
と書ける。（以上(2)）
\end{proof}

\end{document}
